\subsection{Contemporary transactional state and authority boundaries}

The 2026 literature makes the generic ``external state plus guarded write'' idea too crowded to carry a novelty claim by itself. Mullins and Symons give LLM agents an explicit inspectable epistemic state updated by consistency-preserving announcement-like operations \cite{mullins2026}. MemTX treats a memory write as distinct from a committed belief and implements staged validation, provenance-bearing records, commit gating and cascade repair \cite{memtx2026}. PPMF studies provenance laundering under lossy memory consolidation and enforces source-authority non-amplification at later action time \cite{ppmf2026}; TMA-NM goes further by making authority origin-bound and non-malleable under laundering attacks \cite{tmanm2026}. Commit-time authorization independently shows that a witness that was valid earlier in a trajectory need not authorize a later durable effect unless freshness, causal priority, binding and eligibility still hold at commit time \cite{committime2026}. TOKI makes the memory-write contract more explicit still: contradiction resolution is formulated as a bitemporal operator algebra with typed isolation assumptions, provenance-preserving audit rows and soundness obligations for update pipelines \cite{toki2026}. ReSSERAct combines posterior belief with ledgers for uncertainty, temporal validity, replay and control and uses a shielded admissibility gate before action; its formal analysis shows that freshness and contradiction burden are not recoverable from a belief state alone \cite{ozturk2026resseract}. SAVeR independently audits and minimally repairs candidate belief states before commitment under explicit acceptance constraints rather than treating coherent or consensual reasoning as automatically faithful \cite{yuan2026saver}. Arjmandi's ``LLM proposes, Executive disposes'' instrument is closer to the proposal/canonical-state separation used here: a deterministic executive owns committed belief, the LLM files typed proposals, and pre-registered predictions are checked against observations by code; its ablations also illustrate why a mechanism-selective structural result can be meaningful even when end-task efficacy is null \cite{arjmandi2026}.

AuthMem-Bench supplies a controlled empirical test of this lower-level boundary \cite{zhan2026authority}. It holds claim content and the later task fixed while varying source authority, observes authority upgrades in 48 of 49 consolidator--backbone configurations, and finds a mean unauthorized-action rate of 50.3\% when collapsed memories lack authority metadata. In its selected end-to-end pipeline, predicted persistent labels reduce the observed unauthorized-action rate from 16.9\% to 0.0\% without a material decrease in authorized task success, although omission remains recall-limiting. These results make source-governed operational permission a first-class inherited requirement rather than a RAKL novelty claim. They do not establish the different scientific question of whether evidence licenses a representation, mechanism, identification or decision claim.

Structured state itself is likewise established in scientific and long-horizon agent systems. Millison et al. use an explicit state-machine architecture and strongly typed interfaces to constrain an LLM assistant performing physical-science analysis \cite{millison2026}; Chen's State-Aware Runtime separates generation from canonical state, validation, commit, rollback and audit in a long-horizon runtime architecture \cite{chen2026stateaware}. Recent surveys independently organize persistent-agent memory around lifecycle and governance properties rather than retrieval quality alone: Always-On Agents emphasizes authority, scope, mutability, provenance, recoverability and actionability, while the ACL memory survey traces the progression from raw storage to reusable experience and highlights increasing structure in long-term agent memory \cite{ding2026alwayson,luo2026memorysurvey}. These surveys corroborate rather than replace the specific transaction/state mechanisms above, but they make the nearest-work boundary sharper: the scientific contribution cannot be ``we give an LLM a state machine or a persistent memory.''

These systems are parent mechanisms rather than competitors to be routed around. They establish that persistent agent state needs explicit update and authorization boundaries, and their strongest transaction/provenance invariants should be assimilated into any serious implementation of evidence-governed scientific memory. They also narrow the contribution of the present paper. RAKL does not claim to invent staged belief commit, deterministic proposal/commit separation, bitemporal contradiction-update operators, freshness/contradiction ledgers, belief-state self-auditing, provenance preservation, origin-bound or use-specific authority, authority-collapse evaluation, freshness checks at a durable boundary, or structured agent state. The residual question is scientific: how should such commit discipline interact with context-indexed local scientific views, typed representation/mechanism/identification relations, partially ordered multi-axis authority, unresolved contradiction, partial identification, negative scientific history and bounded open-world saturation?

This distinction is testable. A transactional memory system can be perfectly safe about \emph{who may write} while remaining agnostic about whether a predictive certificate can be coerced into a mechanism claim or whether two context-misaligned observations should be called contradictory. Conversely, a scientific relation algebra is incomplete if its canonical updates can be laundered through stale or origin-free memory. The intended RAKL architecture therefore treats transactional commit, typed temporal update algebra, structured runtime state, evidence-state reconciliation and provenance authority as inherited lower-level safeguards and places scientific-state typing above them rather than presenting either layer as a substitute for the other.
